% ARCHIVED 2026-10-13 from _g_appendix.tex: A (symbol table), B and C (soundness, completeness, independence of the
% pullback and pushforward), E (mechanical redundancy, a placeholder). Cut (user): not measured, no practical impact shown.
% Conservation (completeness) survives as one sentence each in _c_method.tex.

\subsection*{A. Symbol Table}

The following table documents the symbols used in this chapter.

\begin{table}[H]
  \setlength{\tabcolsep}{4pt}
  \renewcommand{\arraystretch}{0.92}
  \footnotesize
    \centering
    \begin{tabular}{@{}c | c | l@{}}
      Sym. & Sec. & Definition \\
      \toprule
      $i$ & --- & index typically associated with problem-space \\
      $j$ & --- & index typically associated with feature-space \\
      \midrule
      $z_i$, $z$, $\mathcal{Z}$ &  & problem-space element, object, and space \\
      $z_\varnothing$, $\bar{z}$ &  & null element and augmented problem-space object \\
      $x_j$, $x$, $\mathcal{X}$ &  & feature-space feature, vector, and space \\
      $\sigma$ &  & feature mapping function \\
      $\phi$ &  & classification function \\
      $l$, $\mathcal{L}$ &  & label and space \\
      \midrule
      $b_i$, $b$, $\mathcal{B}$ &  & problem-space importance, explanation, and space \\
      $a_j$, $a$, $\mathcal{A}$ &  & feature-space importance, explanation, and space \\
      \midrule
      $P$ &  & provenance matrix \\
      $U$ &  & unattributable features \\
      $D$ &  & provenance row sums \\
      $D^\dagger$ &  & provenance row-sum pseudo-inverse \\
      $\tau$ &  & pullback operator \\
      $\psi$ &  & pushforward operator \\
      \bottomrule
    \end{tabular}
  \caption{Symbols used in this chapter.}
  \label{tab:XML:Symbol_Table}
\end{table}

\subsection*{B. Characterizing the Pullback}

To clarify the nature of pullback, we formulate three propositions, stated below:

\proposition[\textit{Soundness of Pullback}]{
If for some $i$ we have $P_{i,j} = 0$ for all $j$ with $a_j \neq 0$, then $\tau(a)_i = 0$. \textit{Proof.} $\tau(a)_i = (Pa)_i = \sum_j P_{i,j} a_j$, where every term is zero since $P_{i,j} = 0$ or $a_j = 0$. \hfill$\square$
}

\proposition[\textit{Completeness of Pullback}]{
\label{prop:pullback-completeness}
The initial attribution mass $\sum_j a_j$ is equal to the final mass $\sum_i \tau(a)_i$. \textit{Proof.} $\sum_i \tau(a)_i = \sum_i (Pa)_i = \sum_i \sum_j P_{i,j} a_j = \sum_j a_j \sum_i P_{i,j} = \sum_j a_j$. \hfill$\square$
}

\proposition[\textit{Independence of Pullback}]{
\label{prop:pullback-independence}
Fix $i$. Let $S_i(a) = \{j: a_j \neq 0 \text{ and } P_{i,j} > 0\}$. If $|S_i(a)| \le 1$, then $\tau(a)_i \in \{0, P_{i,j^\star}\, a_{j^\star}\}$ where $j^\star$ is the (unique, if present) index in $S_i(a)$. \textit{Proof.} $\tau(a)_i = \sum_j P_{i,j} a_j$ and by hypothesis at most one summand is nonzero. \hfill$\square$
}

Soundness says pullback never invents new links: an element only receives attribution from features that actually draw on it in the provenance. Completeness says the total amount of attribution is conserved—pulling back just redistributes what you already had, nothing is created or lost. Independence says that when, for a given element, at most one relevant feature touches it (among those with nonzero attribution), the element’s pulled-back value is just that one feature’s score scaled by its provenance weight—no mixing from multiple features. Globally, if the relevant feature supports don’t overlap on any row, pullback produces a collision-free, easy-to-read explanation.

\subsection*{C. Characterizing the Pushforward}

Analogous propositions for soundness, completeness, and independence follow naturally:

\proposition[\textit{Soundness of Pushforward}]{
If for some $j$ we have $P_{i,j}=0$ for all $i$ with $b_i \neq 0$, then $\psi(b)_j = 0$. \textit{Proof.} $\psi(b)_j = \sum_i P_{i,j} D^\dagger_{i,i} b_i$, and every summand is $0$ since $P_{i,j}=0$ or $b_i=0$. \hfill$\square$
}

\proposition[\textit{Completeness of Pushforward}]{
The initial attribution mass $\sum_i b_i$ is equal to the final mass $\sum_j \psi(b)_j$. \textit{Proof.} $\sum_j \psi(b)_j = \sum_j \sum_i P_{i,j} D^\dagger_{i,i} b_i = \sum_i D^\dagger_{i,i} b_i \sum_j P_{i,j} = \sum_i D^\dagger_{i,i} D_{i,i} b_i = \sum_i b_i$ where the last equality uses $D^\dagger_{i,i}D_{i,i}=1$ when $D_{i,i}>0$ and $0$ otherwise, and by Definition~\ref{def:ProblemSpaceExplanation} $b_i=0$ whenever $D_{i,i}=0$. \hfill$\square$
}

\proposition[\textit{Independence of Pushforward}]{
Fix $j$. Let $T_j(b) = \{i: b_i\neq 0 \text{ and } P_{i,j}>0\}$. If $|T_j(b)|\le 1$, then $\psi(b)_j \in \{0, D^\dagger_{i^\star,i^\star} P_{i^\star,j} b_{i^\star}\}$ where $i^\star$ is the (unique, if present) index in $T_j(b)$. \textit{Proof.} $\psi(b)_j = \sum_i P_{i,j} D^\dagger_{i,i} b_i$ and, by hypothesis, at most one summand is nonzero (terms with $D_{i,i} = 0$ vanish). \hfill$\square$
}

Soundness says pushforward only sends attribution to features that an element actually influences in the provenance; if no such link exists, the feature gets zero. Completeness says the total attribution is preserved—each element’s mass is fully redistributed across the features it touches, so the overall sum stays the same. Independence says that when, for a given feature, at most one relevant element feeds into it, the pushed value is just that single element’s score scaled by the corresponding provenance weight—again, no mixing. If, for a given feature $j$, at most one element with $b_i\neq 0$ satisfies $P_{i,j}>0$, pushforward yields a clean, unmixed feature-space explanation.


\subsection*{E. Mechanical Redundancy}
\label{sec:XML:AppRedundancy}

\begin{proposalbox}
Features with identical provenance cannot be told apart after a pullback, so credit split among such features by an explainer recombines exactly. Planned: the proposition, a duplicate-feature sweep, and the agreement of explainers before and after pullback, net of a permutation control.
\end{proposalbox}


